\documentclass[
    language=english,
    doctype=presentation,
    institution=none,
]{../../omnilatex}

\RequirePackage{config/document-settings}

\title{Introduction to OmniLaTeX}
\author{Jane Doe}
\institute{Department of Computer Science, University of Examples}
\date{April 2026}

\begin{document}

%% ── Title Page ────────────────────────────────────

\maketitle

%% ── Outline ───────────────────────────────────────

\begin{presentationframe}
    \slidetitle{Outline}

    \begin{enumerate}
        \item What is OmniLaTeX?
        \item Key Features
        \item Architecture Overview
        \item Block Environments
        \item Getting Started
    \end{enumerate}
\end{presentationframe}

%% ── What is OmniLaTeX? ────────────────────────────

\begin{presentationframe}
    \slidetitle{What is OmniLaTeX?}

    OmniLaTeX is a universal \LaTeX{} template system designed to simplify
    document creation across academic and professional contexts.

    \begin{presentationblock}[Overview]{Core Idea}
        Write once, compile anywhere. OmniLaTeX provides a single template
        that supports 16+ document types with consistent branding and
        professional typesetting.
    \end{presentationblock}

    \vspace{2mm}
    Built on KOMA-Script and compiled with LuaLaTeX for modern font
    and Unicode support.
\end{presentationframe}

%% ── Key Features ──────────────────────────────────

\begin{presentationframe}
    \slidetitle{Key Features}

    \begin{description}
        \item[Modularity] Each document type is a self-contained profile
        \item[Typography] Consistent, professional typesetting out of the box
        \item[Bibliography] Integrated biblatex with automatic citation management
        \item[Code Listings] Syntax highlighting via the minted package
        \item[Multi-language] Full polyglossia support for international documents
        \item[Reproducibility] Docker and Nix flakes for deterministic builds
    \end{description}
\end{presentationframe}

%% ── Section Divider ───────────────────────────────

\presentationSection{Architecture}

%% ── Architecture ──────────────────────────────────

\begin{presentationframe}
    \slidetitle{Layered Architecture}

    The template follows a layered architecture:

    \begin{enumerate}
        \item \textbf{Core Layer} --- Base class, utilities, colors
        \item \textbf{Typography Layer} --- Fonts, math, lists, typesetting
        \item \textbf{Layout Layer} --- Page geometry, floats, boxes, KOMA config
        \item \textbf{Content Layer} --- Document types, institution profiles
        \item \textbf{Reference Layer} --- Hyperlinks, bibliography, glossaries
    \end{enumerate}

    \begin{noteblock}[Tip]{Customization}
        Each layer is independently configurable, allowing users to override
        any aspect of the template without modifying core files.
    \end{noteblock}
\end{presentationframe}

%% ── Block Environments ────────────────────────────

\begin{presentationframe}
    \slidetitle{Block Environments}

    OmniLaTeX presentations provide several styled block environments:

    \begin{definitionblock}[Definition]{Theorem (Euclid)}
        There are infinitely many prime numbers.
    \end{definitionblock}

    \begin{exampleblock}[Example]{Usage}
        Use \texttt{\textbackslash begin\{presentationblock\}\{Title\}}
        for a standard highlighted block.
    \end{exampleblock}

    \begin{alertblock}[Warning]{Important}
        Always compile with LuaLaTeX for full Unicode support.
    \end{alertblock}
\end{presentationframe}

%% ── Document Types ────────────────────────────────

\begin{presentationframe}
    \slidetitle{Supported Document Types}

    OmniLaTeX supports 16+ document types:

    \begin{multicols}{2}
        \begin{itemize}
            \item Thesis / Dissertation
            \item Technical Report
            \item Journal Article
            \item Conference Paper
            \item Book / Manual
            \item CV / Resume
            \item Cover Letter
            \item Standard / Patent
            \item Poster
            \item Presentation
            \item Letter
            \item Dictionary
            \item Inline Paper
        \end{itemize}
    \end{multicols}
\end{presentationframe}

%% ── Getting Started ───────────────────────────────

\begin{presentationframe}
    \slidetitle{Getting Started}

    \begin{presentationblock}[Quick Start]{Minimal Example}
        \ttfamily\small
        \textbackslash documentclass[doctype=presentation]\{omnilatex\}\\
        \textbackslash title\{My Talk\}\\
        \textbackslash author\{Jane Doe\}\\
        \textbackslash begin\{document\}\\
        \quad\textbackslash maketitle\\
        \quad\textbackslash begin\{presentationframe\}\\
        \quad\quad\textbackslash slidetitle\{Hello World\}\\
        \quad\quad Content goes here.\\
        \quad\textbackslash end\{presentationframe\}\\
        \textbackslash end\{document\}
    \end{presentationblock}

    \vspace{2mm}
    Compile with: \texttt{lualatex main.tex}
\end{presentationframe}

%% ── Thank You ─────────────────────────────────────

\begin{presentationframe}
    \vfill
    \begin{center}
        {\Huge\bfseries\color{universityprimary}Thank You!}\par
        \vspace{8mm}
        {\Large Questions?}\par
        \vspace{6mm}
        {\normalsize\texttt{github.com/WyattAu/OmniLaTeX-template}}
    \end{center}
    \vfill
\end{presentationframe}

\end{document}
